\documentclass[border=5pt]{standalone}
\usepackage[T1]{fontenc}
\usepackage{tikz}
\usepackage{luacode}

\definecolor{garnet}{HTML}{73000A}
\definecolor{atlantic}{HTML}{466A9F}
\definecolor{congaree}{HTML}{1F414D}
\definecolor{warmgrey}{HTML}{676156}
\definecolor{bk70}{HTML}{5C5C5C}

\begin{document}
\begin{tikzpicture}

\useasboundingbox (-1,-2.5) rectangle (14,2.5);
\clip (-1,-2.5) rectangle (14,2.5);

\begin{luacode}
-- ════════════════════════════════════════════════════════════
-- CONTROLS
-- ════════════════════════════════════════════════════════════
local h       = 1.4       -- half-aperture height (all lenses same)
local d       = 2.5       -- distance between each element
local n_glass = 1.52
local min_wall = 0.08

-- Element positions (left to right)
local x1 = 2.0                    -- plano-convex
local x_ap = x1 + d               -- aperture
local x2 = x_ap + d               -- biconcave
local x3 = x2 + d                 -- biconvex

-- Lens definitions: {Rl, Rr}
local lens1 = {Rl = 2.5,  Rr = 999}    -- plano-convex
local lens2 = {Rl = -3,   Rr = 3}      -- biconcave
local lens3 = {Rl = 3,    Rr = -3}     -- biconvex

-- Aperture
local ap_gap  = 0.6       -- half-gap (opening radius)

-- Sensor
local sensor_x = x3 + d
local sensor_h = 0.8

-- Rays
local n_rays     = 50
local ray_y_min  = -1.3
local ray_y_max  =  1.3
local ray_width  = 0.15
local ray_opacity = 0.45

-- Scene bounds
local xmin = -1
local xmax = 14
local ymin = -2.5
local ymax = 2.5

-- ════════════════════════════════════════════════════════════
-- ENGINE
-- ════════════════════════════════════════════════════════════

local function auto_t(Rl, Rr, h)
    local aRl, aRr = math.abs(Rl), math.abs(Rr)
    local sl = (aRl >= h) and (aRl - math.sqrt(aRl*aRl - h*h)) or aRl
    local sr = (aRr >= h) and (aRr - math.sqrt(aRr*aRr - h*h)) or aRr
    local dl = (Rl > 0) and -sl or sl
    local dr = (Rr < 0) and -sr or sr
    local t = min_wall - math.min(0, dl + dr)
    if t < min_wall then t = min_wall end
    return t
end

local function isect(ox, oy, dx, dy, cx, cy, r, pick)
    local ex, ey = ox - cx, oy - cy
    local a = dx*dx + dy*dy
    local b = 2*(ex*dx + ey*dy)
    local c = ex*ex + ey*ey - r*r
    local disc = b*b - 4*a*c
    if disc < 0 then return nil end
    local sq = math.sqrt(disc)
    local t1 = (-b - sq)/(2*a)
    local t2 = (-b + sq)/(2*a)
    local t
    if pick == "far" then
        t = (t2 > 1e-4) and t2 or ((t1 > 1e-4) and t1 or nil)
    else
        t = (t1 > 1e-4) and t1 or ((t2 > 1e-4) and t2 or nil)
    end
    if not t then return nil end
    local hx, hy = ox + t*dx, oy + t*dy
    return t, hx, hy, (hx-cx)/r, (hy-cy)/r
end

local function refract(dx, dy, nx, ny, n1, n2)
    local d = dx*nx + dy*ny
    if d > 0 then nx, ny, d = -nx, -ny, -d end
    local r = n1/n2
    local ci = -d
    local s2t = r*r*(1 - ci*ci)
    if s2t > 1 then return nil end
    local ct = math.sqrt(1 - s2t)
    local rx = r*dx + (r*ci - ct)*nx
    local ry = r*dy + (r*ci - ct)*ny
    local l = math.sqrt(rx*rx + ry*ry)
    return rx/l, ry/l
end

local function lens_geo(lx, Rl, Rr, h)
    local t = auto_t(Rl, Rr, h)
    local vl = lx - t/2
    local vr = lx + t/2
    return { vl=vl, vr=vr, t=t,
             cl_x = vl + Rl, cr_x = vr + Rr,
             aRl=math.abs(Rl), aRr=math.abs(Rr) }
end

local function draw_lens(lx, Rl, Rr, h)
    local G = lens_geo(lx, Rl, Rr, h)
    local N = 80
    local lp, rp = {}, {}
    for i = 0, N do
        local y = -h + 2*h*(i/N)
        local ql = math.sqrt(math.max(G.aRl*G.aRl - y*y, 0))
        local qr = math.sqrt(math.max(G.aRr*G.aRr - y*y, 0))
        table.insert(lp, {(Rl > 0) and (G.cl_x - ql) or (G.cl_x + ql), y})
        table.insert(rp, {(Rr < 0) and (G.cr_x + qr) or (G.cr_x - qr), y})
    end
    local p = "\\fill[atlantic!12, opacity=0.85] "
    for i, v in ipairs(lp) do
        p = p .. (i==1 and "" or " -- ") .. string.format("(%.4f,%.4f)", v[1], v[2])
    end
    for i = #rp, 1, -1 do
        p = p .. string.format(" -- (%.4f,%.4f)", rp[i][1], rp[i][2])
    end
    tex.sprint(p .. " -- cycle;")
    for _, arc in ipairs({lp, rp}) do
        local s = "\\draw[atlantic!60, thin] "
        for i, v in ipairs(arc) do
            s = s .. (i==1 and "" or " -- ") .. string.format("(%.4f,%.4f)", v[1], v[2])
        end
        tex.sprint(s .. ";")
    end
    tex.sprint(string.format("\\draw[atlantic!60, thin] (%.4f,%.4f)--(%.4f,%.4f);",
        lp[1][1], lp[1][2], rp[1][1], rp[1][2]))
    tex.sprint(string.format("\\draw[atlantic!60, thin] (%.4f,%.4f)--(%.4f,%.4f);",
        lp[#lp][1], lp[#lp][2], rp[#rp][1], rp[#rp][2]))
end

-- Trace through a single lens element
local function trace_lens(cx, cy, cdx, cdy, lx, Rl, Rr, h)
    local G = lens_geo(lx, Rl, Rr, h)
    local pts = {}
    local pick_l, pick_r
    if cdx >= 0 then
        pick_l = (Rl > 0) and "near" or "far"
        pick_r = (Rr < 0) and "far" or "near"
    else
        pick_l = (Rr < 0) and "near" or "far"
        pick_r = (Rl > 0) and "far" or "near"
    end
    local surf1_cx = (cdx >= 0) and G.cl_x or G.cr_x
    local surf1_r  = (cdx >= 0) and G.aRl or G.aRr
    local surf2_cx = (cdx >= 0) and G.cr_x or G.cl_x
    local surf2_r  = (cdx >= 0) and G.aRr or G.aRl

    local t, hx, hy, nx, ny = isect(cx, cy, cdx, cdy, surf1_cx, 0, surf1_r, pick_l)
    if t and math.abs(hy) <= h then
        table.insert(pts, {hx, hy})
        local ndx, ndy = refract(cdx, cdy, nx, ny, 1.0, n_glass)
        if not ndx then return pts, cdx, cdy, cx, cy end
        cdx, cdy = ndx, ndy
        cx, cy = hx, hy

        t, hx, hy, nx, ny = isect(cx, cy, cdx, cdy, surf2_cx, 0, surf2_r, pick_r)
        if t and math.abs(hy) <= h then
            table.insert(pts, {hx, hy})
            ndx, ndy = refract(cdx, cdy, nx, ny, n_glass, 1.0)
            if not ndx then return pts, cdx, cdy, cx, cy end
            cdx, cdy = ndx, ndy
            cx, cy = hx, hy
        end
    end
    return pts, cdx, cdy, cx, cy
end

-- Full trace: parallel ray from left through all 3 lenses + aperture
local function trace_full(oy)
    local pts = {{xmin, oy}}
    local cx, cy = xmin, oy
    local cdx, cdy = 1, 0

    -- Lens 1: plano-convex
    local lpts, ndx, ndy, nx, ny
    lpts, cdx, cdy, cx, cy = trace_lens(cx, cy, cdx, cdy, x1, lens1.Rl, lens1.Rr, h)
    for _, p in ipairs(lpts) do table.insert(pts, p) end

    -- Aperture: check if ray passes through the gap
    if math.abs(cdx) > 1e-6 then
        local t_ap = (x_ap - cx) / cdx
        if t_ap > 0 then
            local ap_y = cy + t_ap * cdy
            table.insert(pts, {x_ap, ap_y})
            if math.abs(ap_y) > ap_gap then
                -- Blocked by aperture
                return pts, true
            end
            cx, cy = x_ap, ap_y
        end
    end

    -- Lens 2: biconcave
    lpts, cdx, cdy, cx, cy = trace_lens(cx, cy, cdx, cdy, x2, lens2.Rl, lens2.Rr, h)
    for _, p in ipairs(lpts) do table.insert(pts, p) end

    -- Lens 3: biconvex
    lpts, cdx, cdy, cx, cy = trace_lens(cx, cy, cdx, cdy, x3, lens3.Rl, lens3.Rr, h)
    for _, p in ipairs(lpts) do table.insert(pts, p) end

    -- Extend to right edge
    if math.abs(cdx) > 1e-6 then
        local ts = (xmax - cx) / cdx
        if ts > 0 then table.insert(pts, {xmax, cy + ts*cdy}) end
    end
    return pts, false
end

-- ══════════ DRAW ══════════

-- Optical axis
tex.sprint(string.format(
    "\\draw[warmgrey, dashed, thin] (%.2f,0)--(%.2f,0);", xmin, xmax))

-- Lenses
draw_lens(x1, lens1.Rl, lens1.Rr, h)
draw_lens(x2, lens2.Rl, lens2.Rr, h)
draw_lens(x3, lens3.Rl, lens3.Rr, h)

-- Aperture
tex.sprint(string.format(
    "\\draw[bk70, thick] (%.2f,%.2f)--(%.2f,%.2f);",
    x_ap, h + 0.1, x_ap, ap_gap))
tex.sprint(string.format(
    "\\draw[bk70, thick] (%.2f,%.2f)--(%.2f,%.2f);",
    x_ap, -ap_gap, x_ap, -h - 0.1))

-- CMOS sensor
tex.sprint(string.format(
    "\\fill[congaree!80] (%.2f,%.2f) rectangle (%.2f,%.2f);",
    sensor_x, -sensor_h, sensor_x + 0.12, sensor_h))
local npix = 12
for i = 0, npix - 1 do
    local py = -sensor_h + i * (2*sensor_h / npix)
    tex.sprint(string.format(
        "\\draw[congaree!40, very thin] (%.2f,%.4f)--(%.2f,%.4f);",
        sensor_x, py, sensor_x + 0.12, py))
end
tex.sprint(string.format(
    "\\draw[congaree, thick] (%.2f,%.2f) rectangle (%.2f,%.2f);",
    sensor_x, -sensor_h, sensor_x + 0.12, sensor_h))

-- Rays
for i = 0, n_rays - 1 do
    local frac = (n_rays > 1) and (i / (n_rays - 1)) or 0.5
    local y0 = ray_y_min + (ray_y_max - ray_y_min) * frac
    local pts, blocked = trace_full(y0)
    local color = blocked and "warmgrey" or "garnet"
    local opacity = blocked and 0.15 or ray_opacity
    for j = 1, #pts - 1 do
        tex.sprint(string.format(
            "\\draw[%s, line width=%.2fpt, opacity=%.2f] (%.4f,%.4f)--(%.4f,%.4f);",
            color, ray_width, opacity,
            pts[j][1], pts[j][2], pts[j+1][1], pts[j+1][2]))
    end
end

-- Labels
tex.sprint(string.format(
    "\\node[font=\\sffamily\\footnotesize,text=atlantic,anchor=south] at (%.2f,%.2f) {Plano-convex};",
    x1, h + 0.15))
tex.sprint(string.format(
    "\\node[font=\\sffamily\\footnotesize,text=atlantic,anchor=south] at (%.2f,%.2f) {Biconcave};",
    x2, h + 0.15))
tex.sprint(string.format(
    "\\node[font=\\sffamily\\footnotesize,text=atlantic,anchor=south] at (%.2f,%.2f) {Biconvex};",
    x3, h + 0.15))
tex.sprint(string.format(
    "\\node[font=\\sffamily\\footnotesize,text=bk70,anchor=south] at (%.2f,%.2f) {Aperture};",
    x_ap, h + 0.15))
tex.sprint(string.format(
    "\\node[font=\\sffamily\\footnotesize,text=congaree,anchor=south] at (%.2f,%.2f) {CMOS};",
    sensor_x + 0.06, sensor_h + 0.15))

\end{luacode}

\end{tikzpicture}
\end{document}
